\documentclass[12pt]{article}
\usepackage[utf8]{inputenc}
\usepackage[french]{babel}
\usepackage{amsmath}
\usepackage{circuitikz}
\usepackage{pgfplots}
\usepackage{siunitx}

\begin{document}

\section*{Circuit RC Complet}

\subsection*{Schéma du Circuit}
\begin{tikzpicture}[american, font=\sffamily]
    \draw (0,0) to[V, l=$\mathcal{V}_S$] ++(0,4);
    \draw (0,0) to[R, l=$R$] (0,4);
    \draw (0,4) to[C={i>\pgfkeysvalueof{/pgf/minimum width}}, l_=$C$] ++(5,0);
    \draw (5,4) -- (5,0) -- (0,0);
\end{tikzpicture}

\subsection*{Courbe de Charge}
\begin{center}
\begin{tikzpicture}
\begin{axis}[
    xlabel={Temps $t$ (\si{\second})},
    ylabel={$V_C(t)$ (\si{\volt})},
    xmin=0, xmax=0.01,
    ymin=0, ymax=15,
    samples=200,
    grid,
    width=10cm,
    height=6cm
]
\addplot[blue, domain=0:0.01] {15*(1-exp(-x/0.005))};
\addlegendentry{$\mathcal{V}_S(1-e^{-t/RC})$};
\end{axis}
\end{tikzpicture}
\end{center}

\subsection*{Équations}
La tension aux bornes du condensateur :
\[
V_C(t) = \mathcal{V}_S \left(1 - e^{-t/RC}\right)
\]

L'énergie emmagasinée :
\[
E(t) = \frac{1}{2}C(V_C(t))^2 = \frac{\mathcal{V}_S^2 C}{2}\left(1-e^{-t/RC}\right)^2
\]

\end{document}